\documentclass[11pt,aspectratio=169]{beamer}

\usepackage[T1]{fontenc}
\usepackage{lmodern}
\usepackage{amsmath,amssymb,mathtools}
\usepackage{booktabs}
\usepackage{tikz}
\usetikzlibrary{arrows.meta,positioning,calc}
\usepackage{appendixnumberbeamer}

\setbeamertemplate{navigation symbols}{}
\setbeamertemplate{itemize items}[circle]
\setbeamertemplate{itemize subitem}{--}
\setbeamertemplate{footline}[frame number]
\setbeamerfont{framesubtitle}{size=\small}
\setbeamercolor{block title}{fg=black,bg=white!82!black}
\setbeamercolor{block body}{fg=black,bg=white!97!black}
\setbeamercolor{alerted text}{fg=blue!55!black}
\linespread{1.12}

\newcommand{\dd}{\mathrm d}
\newcommand{\1}{\mathbf 1}
\newcommand{\HH}{\mathcal H}
\newcommand{\Mbar}{\bar M}

\title{Intergenerational Housing Allocation and Fertility}
\subtitle{Part I: Compact Analytical Model}
\author{Tommaso De Santo}
\date{June 2026}

\begin{document}

% --------------------------------------------------------------------
\begin{frame}[plain,noframenumbering]
\titlepage
\vspace{-0.8em}
\begin{center}
\small
Question: when does limited access to family-capable housing reduce fertility, and when is policy effective?
\end{center}
\end{frame}

% --------------------------------------------------------------------
\begin{frame}{1. Primitives and Timing}
\small
\textbf{Households.}
A mass \(\Mbar\) of potential young households has type \(i\), drawn from \(G_Y\). A type carries income \(y_i\), pledgeable wealth \(a_i\), inheritance exposure \(B_i\), and child consumption cost \(\chi_i\). Old incumbent owners are indexed by \(j\), drawn from \(F_O\).

\medskip
\textbf{Timing for young households.}
\begin{enumerate}
  \item Decide whether to enter the modeled housing market.
  \item Conditional on entry, solve renter and owner branch problems.
  \item Choose tenure \(o_i\in\{R,O\}\), housing \(h_i\), residual adult consumption \(c_i\), and fertility \(n_i\).
\end{enumerate}

\medskip
\textbf{Child inputs.}
\[
c_i=C_i-\chi_i n_i,\qquad s_i=h_i-\kappa n_i,
\]
\[
u_i^Y(c_i,h_i,n_i)
=
\log c_i+\alpha\log(h_i-\kappa n_i)+\beta\log n_i.
\]
\(\alpha\) is the adult housing-services weight, \(\beta\) is the fertility weight, and \(\kappa\) is housing services per child.
\end{frame}

% --------------------------------------------------------------------
\begin{frame}{2. Housing Segments, Prices, and Supply}
\footnotesize
\textbf{Two housing segments.}
\[
\HH^R=\{h:0<h\le \bar h^R\},\qquad
\HH^O=\{h:0<h\le \bar h^O\},\qquad
0<\bar h^R<\bar h^O.
\]
\(\bar h^m\) is an individual occupancy cap. Aggregate stock is \(H^m\).

\smallskip
\textbf{Stocks and construction.}
\[
H^m=\bar H^m+I^m,\qquad
q^m=\Phi_I^m(I^m;Z^m),\qquad m\in\{R,O\}.
\]
\(\bar H^m\) is inherited stock, \(I^m\) is net construction, \(\Phi^m\) is the flow resource cost, and \(Z^m\) shifts construction ease.

\smallskip
\textbf{Owner stock accounting.}
\[
\int H_j^0\,\dd F_O(j)+H_{\mathrm{pass}}^O=\bar H^O.
\]
Old owners enter with endowments \(H_j^0\); remaining owner stock is held by a passive account.

\smallskip
\textbf{Owner price.}
\[
P^O=\frac{q^O}{r+\delta^O+\tau^p}.
\]
\(q^O\) is the flow user cost; \(P^O\) is the asset price used in down-payment constraints. \(r\) is the interest rate, \(\delta^O\) depreciation, and \(\tau^p\) the property-tax rate.
\end{frame}

% --------------------------------------------------------------------
\begin{frame}{3. Young Household Problem}
\small
\textbf{Pledgeable resources.}
\[
x=1-\tau^b,\qquad A_i(x)=a_i+xB_i+T_i^b(x).
\]
\(A_i(x)\) enters purchase constraints, not the flow budget.

\medskip
\textbf{Renter branch.}
\[
W_i^R(q^R)=\max_{c,h,n} u_i^Y(c,h,n)
\]
\[
c+q^R h+\chi_i n=y_i,\qquad h\le \bar h^R,\qquad h-\kappa n>0.
\]

\medskip
\textbf{Owner branch.}
\[
W_i^O(q^O,P^O,x)=\max_{c,h,n} u_i^Y(c,h,n)
\]
\[
c+q^O h+\chi_i n=y_i,\quad
h\le \bar h^O,\quad
(1-\phi)P^O h\le A_i(x),\quad
q^O h\le \psi y_i.
\]
\(\phi\) is the financed share; \(\psi\) is the payment-to-income limit.
\end{frame}

% --------------------------------------------------------------------
\begin{frame}{4. Tenure, Entry, and Old Owners}
\footnotesize
\textbf{Tenure choice conditional on entry.}
\[
W_i^Y(q^R,q^O,P^O,x)=\max\{W_i^R(q^R),W_i^O(q^O,P^O,x)\}.
\]

\smallskip
\textbf{Logit entry.}
\[
\pi_i^E
=
\frac{\exp\{W_i^Y/\kappa_E\}}
{\exp\{W_i^Y/\kappa_E\}+\exp\{\bar W^E/\kappa_E\}}.
\]
\(\bar W^E\) is the outside option and \(\kappa_E\) is the entry taste-shock scale.

\smallskip
\textbf{Old incumbent owners.}
Old owner \(j\) enters with owner housing endowment \(H_j^0\). If she keeps the current house, she pays tax on a protected assessment. The per-unit tax discount from staying put is
\[
L_j^\tau=\tau^pP^O-\tilde \tau_j\tilde P_j\ge0.
\]
Here \(\tau^pP^O\) is the tax cost at current market value and \(\tilde\tau_j\tilde P_j\) is the protected tax cost. Old owner \(j\) chooses how much current housing to keep, consumption, and bequest expenditure:
\[
\max_{c,h,e} u_j^O(c,h)+\mathcal B_j(e)
\quad\text{s.t.}\quad
c+e+(q^O-L_j^\tau)h=m_j+q^O H_j^0,\qquad 0\le h\le H_j^0.
\]
\(m_j\) is nonhousing resources. She feels carrying cost \(q^O-L_j^\tau\):
\[
MV_j^O=q^O-L_j^\tau.
\]
\end{frame}

% --------------------------------------------------------------------
\begin{frame}{5. Competitive Equilibrium}
\footnotesize
Given policies and stocks, an equilibrium is
\[
\left(q^R,q^O,P^O,I^R,I^O,\{\pi_i^E,o_i,c_i,h_i,n_i\}_i,\{c_j^O,h_j^O,e_j\}_j\right)
\]
such that:

\smallskip
\textbf{1. Household optimization.}
Young households solve renter and owner branches, choose tenure by \(W_i^Y=\max\{W_i^R,W_i^O\}\), and enter with probability \(\pi_i^E\). Old owners choose consumption, bequests, and how much current housing to keep.

\smallskip
\textbf{2. Construction and asset pricing.}
\[
q^R=\Phi_I^R(I^R;Z^R),\qquad
q^O=\Phi_I^O(I^O;Z^O),\qquad
P^O=\frac{q^O}{r+\delta^O+\tau^p}.
\]

\smallskip
\textbf{3. Segment clearing.}
\[
\Mbar\int \pi_i^E\1\{o_i=R\}h_i^R\,\dd G_Y(i)=\bar H^R+I^R,
\]
\[
\Mbar\int \pi_i^E\1\{o_i=O\}h_i^O\,\dd G_Y(i)+\int h_j^O\,\dd F_O(j)=\bar H^O+I^O.
\]

\smallskip
\textbf{4. Accounts.}
Fiscal revenues, passive housing income, and policy-account costs are rebated through the transfer rules already embedded in \(A_i(x)\), \(y_i\), and old resources.
\end{frame}

% --------------------------------------------------------------------
\begin{frame}{6. Planner Problem}
\footnotesize
\textbf{Choices.}
The planner chooses tenure assignments \(o_i^P\), young allocations \((c_i,h_i,n_i)\), entry values \(\pi_i^P\), old allocations \((c_j^O,h_j^O,e_j)\), construction \((I^R,I^O)\), and policy instruments.

\smallskip
\textbf{Segment feasibility.}
\[
\Mbar\int \pi_i^P\1\{o_i^P=R\}h_i\,\dd G_Y(i)\le \bar H^R+I^R,
\]
\[
\Mbar\int \pi_i^P\1\{o_i^P=O\}h_i\,\dd G_Y(i)+\int h_j^O\,\dd F_O(j)\le \bar H^O+I^O.
\]

\smallskip
\textbf{Goods feasibility.}
\[
\Mbar\int \pi_i^P(c_i+\chi_i n_i-y_i)\,\dd G_Y(i)
+\int(c_j^O+e_j-m_j)\,\dd F_O(j)
+\Phi^R+\Phi^O+\Xi^\Pi\le0.
\]

\smallskip
\textbf{Implementation.}
Menus \(h_i\in\HH^m\) are real constraints. Down-payment and payment-to-income constraints remain unless the planner uses feasible instruments whose resource costs enter \(\Xi^\Pi\).
\end{frame}

% --------------------------------------------------------------------
\begin{frame}{7. Planner Marginal Conditions}
\footnotesize
\textbf{Segment shadow prices.}
Let \(\Lambda Q^R\) and \(\Lambda Q^O\) be the multipliers on rental and owner segment feasibility; \(\Lambda\) is the goods multiplier.

\smallskip
\textbf{Interior allocation rules away from real menu corners.}
\[
\frac{u_{h,i}^Y}{u_{c,i}^Y}=Q^R
\quad\text{for renter assignments},
\qquad
\frac{u_{h,i}^Y}{u_{c,i}^Y}=Q^O
\quad\text{for young owner assignments},
\]
\[
\frac{u_{h,j}^O}{u_{c,j}^O}=Q^O
\quad\text{for old owner housing},
\qquad
Q^m=\Phi_I^m(I^m;Z^m).
\]

\smallskip
\textbf{Fertility rule.}
\[
\frac{\beta c_i}{n_i}+\nu_i^g=\chi_i+\kappa Q^{o_i^P}.
\]
\(\nu_i^g\) is the goods-equivalent social value of births beyond parental utility; in the fertility-neutral planner, \(\nu_i^g=0\).

\smallskip
\textbf{Connection to CE prices.}
At an interior competitive allocation with construction and segment clearing,
\[
Q^R=q^R,\qquad Q^O=q^O.
\]
\end{frame}

% --------------------------------------------------------------------
\begin{frame}{8. Where the CE Differs from the Planner}
\footnotesize
\textbf{Young owner financial constraints.}
If down-payment or payment-to-income constraints bind, the competitive owner margin is
\[
MV_i^O=q^O+\zeta_i^{O,F},
\qquad
\zeta_i^{O,F}
=
\frac{\mu_i}{\lambda_i^O}(1-\phi)P^O
+
\frac{\eta_i}{\lambda_i^O}q^O.
\]
\[
\underbrace{MV_i^O}_{\text{young private value}}
-
\underbrace{Q^O}_{\text{planner scarcity cost}}
=
\underbrace{\zeta_i^{O,F}}_{\text{financial access gap}}
\qquad (Q^O=q^O).
\]
\(\lambda_i^O\) is the owner budget multiplier; \(\mu_i\) and \(\eta_i\) are down-payment and payment-to-income multipliers.

\smallskip
\textbf{Old owners: incumbent tax discount.}
\[
MV_j^O=q^O-L_j^\tau,
\qquad
\underbrace{Q^O}_{\text{planner scarcity cost}}
-
\underbrace{MV_j^O}_{\text{old private value}}
=
\underbrace{L_j^\tau}_{\text{tax discount from staying put}}.
\]

\smallskip
\textbf{Efficiency statement.}
With zero marginal resource cost of relaxing DP/PTI, positive \(\zeta_i^{O,F}\) or \(L_j^\tau\) violates the planner's marginal conditions. With positive instrument cost, compare the gap to that cost.
\end{frame}

% --------------------------------------------------------------------
\begin{frame}{9. Fertility Effect of a Family-Space Cap}
\footnotesize
\textbf{Effective cap in each branch.}
\[
H_i^R=\bar h^R,\qquad
H_i^O=\min\left\{
\bar h^O,\,
\frac{A_i(x)\rho_p}{(1-\phi)q^O},\,
\frac{\psi y_i}{q^O}
\right\},
\qquad
\rho_p=r+\delta^O+\tau^p.
\]

\smallskip
\textbf{Within-branch response.}
For a fixed tenure branch with user cost \(q\), effective cap \(H\), and residual consumption and housing \((c,s)\),
\[
\Psi_i^m
\equiv
\frac{\partial n_i^m}{\partial H_i^m}
=
\frac{\alpha\kappa/s^2-\chi_i q^m/c^2}
{\beta/n^2+\chi_i^2/c^2+\alpha\kappa^2/s^2}.
\]

\smallskip
\textbf{Sign.}
If the child input bundle is not more consumption-intensive than the adult residual bundle,
\[
\chi_i\le \frac{\kappa q^m}{\alpha},
\]
then relaxing a strictly binding family-housing cap raises fertility within the fixed branch.
\end{frame}

% --------------------------------------------------------------------
\begin{frame}{10. Local Policy Decomposition}
\scriptsize
\textbf{Objects.}
For policy \(z\), let \(P_i^m=\partial H_i^m/\partial z\), \(\Theta_i^m=\lambda_i^m\zeta_i^m\), and \(\mathcal B_m\) be types whose chosen branch \(m\) has a binding effective cap.

\smallskip
\textbf{Fixed-price, fixed-tenure, fixed-active-set formula.}
\[
\begin{aligned}
\left.\frac{\dd N}{\dd z}\right|_{q,o,\mathcal A}
&=
\Mbar\sum_m\int_{\mathcal B_m}
\underbrace{\pi_i^E}_{\substack{\text{entered}\\ \text{mass}}}
\underbrace{\Psi_i^m}_{\substack{\text{fertility}\\ \text{response}}}
\underbrace{P_i^m}_{\substack{\text{policy pass-through}\\ \text{to }H_i^m}}
\,\dd G_Y(i)
\\
&\quad+
\Mbar\sum_m\int_{\mathcal B_m}
\underbrace{(n_i^m-\bar n_i^E)}_{\substack{\text{births relative}\\ \text{to outside option}}}
\underbrace{\frac{\pi_i^E(1-\pi_i^E)}{\kappa_E}}_{\substack{\text{entry}\\ \text{sensitivity}}}
\underbrace{\Theta_i^m}_{\substack{\text{utility gain}\\ \text{from }H_i^m}}
\underbrace{P_i^m}_{\substack{\text{policy pass-through}\\ \text{to }H_i^m}}
\,\dd G_Y(i).
\end{aligned}
\]

\smallskip
\textbf{Reading it.}
The first line is the intensive fertility response of already-entered constrained households. The second line is the entry response. Tenure switching, price feedback, construction responses, and resource costs \(\Xi^\Pi\) are outside this local formula.
\end{frame}

% --------------------------------------------------------------------
\appendix

\begin{frame}{Appendix: Branch Derivative}
\small
If \(0<H<h^u\), then \(h(H)=H\) and
\[
c(H)=y-qH-\chi n(H),
\qquad
s(H)=H-\kappa n(H).
\]
The constrained fertility condition is
\[
F(n,H)\equiv
\frac{\beta}{n}
-
\frac{\chi}{c(H)}
-
\frac{\alpha\kappa}{s(H)}
=0.
\]
\[
F_n
=
-
\left[
\frac{\beta}{n^2}
+
\frac{\chi^2}{c^2}
+
\frac{\alpha\kappa^2}{s^2}
\right]<0,
\qquad
F_H
=
\frac{\alpha\kappa}{s^2}
-
\frac{\chi q}{c^2}.
\]
\[
\frac{\dd n}{\dd H}
=
-\frac{F_H}{F_n}
=
\frac{
\frac{\alpha\kappa}{s^2}
-
\frac{\chi q}{c^2}
}{
\frac{\beta}{n^2}
+
\frac{\chi^2}{c^2}
+
\frac{\alpha\kappa^2}{s^2}
}.
\]
\end{frame}

\begin{frame}{Appendix: Equal Scaling and Access Gaps}
\small
At a binding cap, the household marginal value of housing is
\[
\frac{u_h}{u_c}
=
\frac{\alpha c_i^m}{s_i^m}
=q^m+\zeta_i^m.
\]
Under exact equal scaling, \(\chi_i=\kappa q^m/\alpha\). The numerator of \(\Psi_i^m\) is
\[
\frac{\alpha\kappa}{(s_i^m)^2}
-
\frac{\chi_i q^m}{(c_i^m)^2}
=
\frac{\kappa}{\alpha}
\left(
\frac{\alpha}{s_i^m}-\frac{q^m}{c_i^m}
\right)
\left(
\frac{\alpha}{s_i^m}+\frac{q^m}{c_i^m}
\right).
\]
Since
\[
\frac{\alpha}{s_i^m}-\frac{q^m}{c_i^m}
=
\frac{\zeta_i^m}{c_i^m},
\]
\[
\Psi_i^m
=
\frac{\kappa}{\alpha}
\frac{\zeta_i^m}{c_i^m}
\Omega_i^m,
\qquad
\Omega_i^m>0.
\]
\end{frame}

\begin{frame}{Appendix: Planner Comparison}
\small
At an interior competitive allocation with construction,
\[
Q^R=q^R=\Phi_I^R(I^R;Z^R),
\qquad
Q^O=q^O=\Phi_I^O(I^O;Z^O).
\]
Planner housing rules require
\[
MV_i^R=q^R,
\qquad
MV_i^O=q^O,
\qquad
MV_j^O=q^O.
\]
Competitive marginal values are
\[
MV_i^R=q^R+\zeta_i^R,
\qquad
MV_i^O=q^O+\zeta_i^O,
\qquad
MV_j^O=q^O-L_j^\tau.
\]
\medskip
When real menu corners are common to households and planner, the implementation-relevant young gap is the financial component:
\[
\zeta_i^{O,F}
=
\frac{\mu_i}{\lambda_i^O}(1-\phi)P^O
+
\frac{\eta_i}{\lambda_i^O}q^O.
\]
If financial constraints are relaxable at zero marginal real resource cost, \(\zeta_i^{O,F}>0\) violates the planner's marginal condition.
\end{frame}

\end{document}
